% Source: https://artofproblemsolving.com/community/c3h1250950_geometry_and_circles
% Stable identity: aops-topic-1250950-q01
The external tangents meet at a point $P$ on the line of centres, and this line bisects the $60^\circ$ angle, so each tangent makes $30^\circ$ with it.

\begin{asy}
usepackage("amsmath");
size(6cm);
defaultpen(fontsize(10pt));
pair P=(0,0); pair O1=(30,0),O2=(190,0);
draw(circle(O1,15)); draw(circle(O2,95));
draw(P--(230*dir(30)),gray); draw(P--(230*dir(-30)),gray); draw(P--O2,dashed);
pair T1=O1+15*dir(120), T2=O2+95*dir(120); draw(O1--T1,red); draw(O2--T2,red);
dot(P);dot(O1);dot(O2); label("$P$",P,W); label("$O_1$",O1,S); label("$O_2$",O2,S);
label("$15$",(O1+T1)/2,W); label("$95$",(O2+T2)/2,W);
label("\mbox{}",(min(currentpicture,user=true)+max(currentpicture,user=true))/2,black); // keeps the text after the figure black when LaTeX typesets the labels inline
\end{asy}

The radius to a point of tangency is perpendicular to the tangent, so for each circle $\sin30^\circ=\dfrac{r}{PO}$, i.e. $PO=2r$. Thus $PO_1=30$ and $PO_2=190$, and the centres are $190-30=\boxed{160}$ apart.
